\section{方法论沉淀：可复用的训练决策框架}

结合整场内容，可以提炼一个可直接用于团队执行的决策框架：先定义目标与约束，再搭建 grader 栈，随后构建数据闭环，最后做系统效率优化。顺序错了，常常会产生高成本返工。

\begin{table}[H]
\centering
\begin{tabular}{p{2.0cm}p{4.2cm}p{4.6cm}}
\toprule
\textbf{阶段} & \textbf{关键动作} & \textbf{输出物} \\
\midrule
问题定义 & 任务边界、成功标准、失败清单 & 目标规范文档、风险清单 \\
评价设计 & verifiers + rubric + model graders & grader DAG、冲突矩阵、权重策略 \\
数据构建 & 真实样本、合成样本、回流样本混训 & 数据配方、采样策略、版本控制 \\
训练执行 & 异步 rollout、RL 更新、监控回路 & 可复现训练作业、对比实验记录 \\
部署反馈 & 线上监控、失败回放、再训练触发 & 迭代 backlog、数据增量计划 \\
\bottomrule
\end{tabular}
\caption{从课程抽象出的 Agent 训练执行框架}
\end{table}

\begin{importantbox}{实践准则：先把 ``可测'' 做好，再追求 ``极限''}
在模型还不稳定时，优先建立可重复评估和高吞吐迭代机制，比追逐单次 SOTA 分数更有效。因为后者不可复现时几乎无法转化为产品价值。
\end{importantbox}

另一个实践启发是 ``suboptimal but faster'' 的基线哲学。先用更便宜配置跑更多假设验证，确认有效后再上高成本最优配置，能显著提升单位 GPU 的创新产出。

\begin{knowledgebox}{课程对研究生和工程师的共同价值}
研究生可从中学习 ``如何把想法转化为可验证实验''；工程师可学习 ``如何把实验结果转化为产品行为改进''。两者之间的桥梁正是 grader 与数据闭环。
\end{knowledgebox}

\subsection{本章小结}
这节课的可复用价值在于方法论：用系统工程视角管理 Agent 训练，而不是把它视为单点算法优化。

